%======================================================================
%  notation.tex -- one place for the notation of the whole book
%
%  scalars ............ italic                 u, \theta, \lambda
%  vectors ............ bold italic            \vect{u}, \vu, \vn
%  2nd-order tensors .. bold italic            \tens{\sigma}, \sig, \eps
%  4th-order tensors .. blackboard bold        \tensf{C}, \bbC, \bbS, \bbI
%  matrices (discrete)  bold upright           \mat{K}, \mat{B}
%  column of unknowns . bold upright           \mat{u}  (= [u] in the slides)
%======================================================================
\newcommand{\vect}[1]{\bm{#1}}
\newcommand{\tens}[1]{\bm{#1}}
\newcommand{\tensf}[1]{\mathbb{#1}}
\newcommand{\mat}[1]{\mathbf{#1}}

% sets and spaces
\newcommand{\R}{\mathbb{R}}
\newcommand{\Vsp}{\mathcal{V}}          % functional space of the variational problem
\newcommand{\Kad}{\mathcal{K}}          % kinematically admissible fields
\newcommand{\Sad}{\mathcal{S}}          % statically admissible fields

% vectors
\newcommand{\vu}{\vect{u}}
\newcommand{\vv}{\vect{v}}
\newcommand{\vw}{\vect{w}}
\newcommand{\vx}{\vect{x}}
\newcommand{\vy}{\vect{y}}
\newcommand{\vn}{\vect{n}}
\newcommand{\vt}{\vect{t}}
\newcommand{\vf}{\vect{f}}
\newcommand{\vb}{\vect{b}}
\newcommand{\ve}{\vect{e}}
\newcommand{\vd}{\vect{d}}
\newcommand{\vs}{\vect{s}}
\newcommand{\vq}{\vect{q}}
\newcommand{\vp}{\vect{p}}
\newcommand{\vF}{\vect{F}}
\newcommand{\vE}{\vect{E}}
\newcommand{\vm}{\vect{m}}
\newcommand{\vxi}{\vect{\xi}}
\newcommand{\vnu}{\vect{\nu}}
\newcommand{\vlambda}{\vect{\lambda}}
\newcommand{\vzero}{\vect{0}}

% second-order tensors
\newcommand{\sig}{\tens{\sigma}}
\newcommand{\eps}{\tens{\varepsilon}}
\newcommand{\tI}{\tens{I}}               % identity
\newcommand{\tk}{\tens{k}}               % conductivity
\newcommand{\tA}{\tens{A}}               % coefficient field A(x)
\newcommand{\tM}{\tens{M}}               % polarization tensor (Ch. 3)
\newcommand{\tmu}{\tens{\mu}}            % moment tensor of the opening (Ch. 3)
\newcommand{\talpha}{\tens{\alpha}}      % thermal expansion

% fourth-order tensors
\newcommand{\bbC}{\tensf{C}}             % elastic moduli (Hooke)
\newcommand{\bbS}{\tensf{S}}             % compliances, S = C^{-1}
\newcommand{\bbI}{\tensf{I}}             % symmetric identity

% operators
\DeclareMathOperator{\dive}{div}
\DeclareMathOperator{\tr}{tr}
\DeclareMathOperator{\dev}{dev}
\DeclareMathOperator{\sym}{sym}
\DeclareMathOperator{\diag}{diag}
\DeclareMathOperator*{\argmin}{arg\,min}
\DeclareMathOperator*{\argmax}{arg\,max}
\newcommand{\T}{^{\!\top}}                % transpose
\newcommand{\dd}{\,\mathrm{d}}            % integration measure
\newcommand{\jump}[1]{[\![#1]\!]}
\newcommand{\norm}[1]{\lVert #1\rVert}
\newcommand{\scal}[2]{\langle #1,#2\rangle}
\newcommand{\dpr}{\!:\!}                   % double contraction (compact)

% boundary data (superscript D = "datum")
\newcommand{\uD}{u^{D}}
\newcommand{\pD}{\varphi^{D}}
\newcommand{\vuD}{\vu^{D}}
\newcommand{\vtD}{\vt^{D}}
\newcommand{\vfD}{\vf^{D}}

% Chapter 3 (reciprocity gap)
\newcommand{\RG}{\mathcal{R}}
\newcommand{\OG}{\Omega_\Gamma}
